\chapter{Development environment settings}
\label{ch:setup}

In this chapter, we set up the necessary development environment prior to full-scale LLM implementation. Explore Python, PyTorch, and the book's code repository structure, and run your first ``Hello, LLM'' program.

\section{tools needed}

\begin{center}
\begin{tabularx}{\textwidth}{lXl}
\toprule
\textbf{equipment}&\textbf{use}&\textbf{Recommended version}\\
\midrule
Python & implementation language & 3.10 or higher \\
PyTorch & Deep Learning Framework & 2.0 or higher \\
NumPy & Numerical Computing & 1.24 and higher \\
pytest & unit test & 7.0+ \\
git & version control & latest \\
\bottomrule
\end{tabularx}
\end{center}

GPU is optional. All code in this book is written to run on the CPU. However, for large-scale learning, CUDA-capable GPUs are strongly recommended. When discussing distributed learning in Volume 2, a multi-GPU environment is assumed.

\section{PyTorch installation}

PyTorch follows the installation guide on the official website. Choose the one that suits your environment between the CPU-only version and the CUDA version.

\begin{lstlisting}[style=bashstyle]
# CPU Only (all examples in this book are CPUoperation)
pip install torch --index-url https://download.pytorch.org/whl/cpu

# CUDA 12.1 Support (GPU when used)
pip install torch --index-url https://download.pytorch.org/whl/cu121
\end{lstlisting}

To check installation, do the following:

\begin{lstlisting}[language=Python]
import torch
print(torch.__version__)         # 2.x.x
print(torch.cuda.is_available()) # GPU Availability
\end{lstlisting}

\section{repository structure}

The code in this book has the following structure. The modules covered in each chapter are separated into separate files, so you can use only the parts you need.

\begin{lstlisting}[style=bashstyle]
make-llm-basic/
  src/makellm/
    __init__.py
    tokenizer/              # Chapter 3: tokenizer
      base.py
      bpe.py
      unigram.py
    model/                  # 4-Chapter 6: model structure
      embedding.py
      attention.py
      transformer_block.py
      gpt.py
    training/               # Chapter 7: learning
      dataset.py
      loss.py
      optimizers.py
      trainer.py
    inference/              # Chapter 8: inference
      sampler.py
      generate.py
    utils/                  # common utility
      config.py
      seed.py
      logging.py
  tests/                    # pytest unit testing
  book/                     # of this book LaTeX sauce
\end{lstlisting}

\section{first example: Dealing with tensors}

Before full-scale implementation, let’s review the basics of PyTorch tensors. Tensors operate almost identically to NumPy arrays, except that they support GPU operations and automatic differentiation.

\begin{lstlisting}[language=Python]
import torch

# Tensor creation
x = torch.randn(2, 3)         # Sampling from a standard normal distribution
print(x.shape)                # torch.Size([2, 3])
print(x.dtype)                # torch.float32

# calculation
y = x @ x.T                   # matrix multiplication: [2, 3] @ [3, 2] -> [2, 2]
z = x.sum()                   # Sum of all elements (scalar)

# automatic differentiation
w = torch.randn(3, requires_grad=True)
loss = (w * w).sum()          # L2 norm^2
loss.backward()               # backward
print(w.grad)                 # d(loss)/d(w) = 2*w
\end{lstlisting}

\begin{infobox}[Autograd]
One of PyTorch's most powerful features. All operations performed on tensors are internally recorded as a ``computation graph'', and gradients for all parameters are automatically calculated\code{loss.backward()}once. There is no need to derive the differential equation ourselves. All models in this book rely on autograd.
\end{infobox}

\section{Ensure reproducibility}

Deep learning experiments have a lot of randomness, so reproducibility is important. If you fix the seed, you can run the same code multiple times and get the same results.

\begin{lstlisting}[language=Python]
from makellm.utils import set_seed
set_seed(42)  # Python, NumPy, PyTorch Set up your seeds in one step
\end{lstlisting}

This function calls Python's\code{random}, NumPy's\code{numpy.random}, and PyTorch's\code{torch.manual\_seed}. If CUDA exists, the CUDA seed is also set.

\section{Run unit tests}

The code for each chapter is provided with pytest-based unit tests. After modifying the code, be sure to run tests to make sure there are no regressions.

\begin{lstlisting}[style=bashstyle]
cd make-llm-basic
pytest tests/ -v
\end{lstlisting}

\begin{lstlisting}[style=bashstyle, basicstyle=\ttfamily\small]
========================= 48 passed in 2.47s =========================
\end{lstlisting}

If all 48 tests pass, the environment setup is complete.

\section{summation}

\begin{itemize}
\item Installed PyTorch 2.x and pytest.
\item I understand the repository structure: Code is separated by module in\code{src/makellm/}.
\item We reviewed the basics of tensor operations and autograd.
\item Use\code{set\_seed()}for reproducibility.
\item You can run unit tests with \code{pytest tests/}.
\end{itemize}

In the next chapter, we will implement the tokenizer in earnest. This first conversion of text to numbers is the starting point for all LLMs.
